\documentclass[11pt,oneside]{article}

\pdfminorversion=7
\usepackage{QuizStyle}

\hypersetup{
    pdftitle={20261015 Quiz for MATH261 Discrete Mathematics},
    pdfauthor={Jungwoo Kim},
    pdfsubject={Quiz for MATH261 Discrete Mathematics},
    pdfcreator={LaTeX}
}

\begin{document}

\quiztitle{20261015 Quiz}


\begin{problem}{1}
Let $m,n\in\N$.
\begin{problemsubparts}
    \item Let $k_1,\ldots,k_m\in\N$ with $k_1+\cdots+k_m=n$. Prove that
    \[
        \binom{n}{k_1,k_2,\ldots,k_m}
        =\sum_{j=1}^{m}
        \binom{n-1}{k_1,\ldots,k_j-1,\ldots,k_m}.
    \]

    \item Prove the multinomial theorem by induction on $n$: for arbitrary $x_1,\ldots,x_m\in\R$,
    \[
        (x_1+\cdots+x_m)^n
        =\sum_{\substack{k_1+\cdots+k_m=n\\k_1,\ldots,k_m\in\Nzero}}
        \binom{n}{k_1,\ldots,k_m}
        x_1^{k_1}\cdots x_m^{k_m}.
    \]
\end{problemsubparts}
\end{problem}


\begin{solution}
\begin{problemsubparts}
    \item Count the strings of length $n$ containing exactly $k_j$ copies of the symbol $j$ for each $j\in[m]$. The total number of such strings is
    \[
        \binom{n}{k_1,k_2,\ldots,k_m}.
    \]
    Classify these strings according to their last symbol. If the last symbol is $j$, then the first $n-1$ positions contain $k_j-1$ copies of $j$ and $k_i$ copies of each $i\neq j$, giving
    \[
        \binom{n-1}{k_1,\ldots,k_j-1,\ldots,k_m}
    \]
    strings. The $m$ cases are disjoint and exhaustive, so summing them gives the desired identity.

    \item We use induction on $n$. For $n=1$,
    \[
        x_1+\cdots+x_m
        =\sum_{j=1}^{m}x_j,
    \]
    which is the multinomial formula in the base case.

    Assume the formula holds for some $n\in\N$. Multiplying the induction hypothesis by $x_1+\cdots+x_m$, consider a monomial
    \[
        x_1^{\ell_1}\cdots x_m^{\ell_m},
        \qquad
        \ell_1+\cdots+\ell_m=n+1.
    \]
    Its coefficient is
    \[
        \sum_{\substack{1\leq j\leq m\\\ell_j>0}}
        \binom{n}{\ell_1,\ldots,\ell_j-1,\ldots,\ell_m}.
    \]
    Using the factorial formula for multinomial coefficients,
    \[
    \begin{aligned}
        \sum_{\substack{1\leq j\leq m\\\ell_j>0}}
        \binom{n}{\ell_1,\ldots,\ell_j-1,\ldots,\ell_m}
        &=\frac{n!}{\ell_1!\cdots\ell_m!}
          \sum_{j=1}^{m}\ell_j\\
        &=\frac{(n+1)!}{\ell_1!\cdots\ell_m!}\\
        &=\binom{n+1}{\ell_1,\ldots,\ell_m}.
    \end{aligned}
    \]
    Therefore the multinomial formula holds for $n+1$, and hence for every $n\in\N$ by mathematical induction.
\end{problemsubparts}
\end{solution}

\begin{problem}{2}
There are $10$ couples, each with one child. The $10$ children are lost in a dark cave, and each couple retrieves exactly one child, so that every child is retrieved once. How many assignments pair at least one child with their own parents?
\end{problem}

\begin{solution}
Label the couples and their children by $[10]$, with child $i$ belonging to couple $i$. An assignment of the children to the couples is a permutation of $[10]$. For $i\in[10]$, let $A_i$ be the set of assignments in which couple $i$ retrieves child $i$.

If $k$ specified couples are correctly paired, then the remaining $10-k$ children may be assigned arbitrarily to the remaining couples, giving $(10-k)!$ assignments. Hence, by inclusion-exclusion,
\[
\begin{aligned}
    |A_1\cup\cdots\cup A_{10}|
    &=\sum_{k=1}^{10}(-1)^{k+1}\binom{10}{k}(10-k)!\\
    &=10!\sum_{k=1}^{10}\frac{(-1)^{k+1}}{k!}\\
    &=2{,}293{,}839.
\end{aligned}
\]
\end{solution}

\begin{problem}{3}
Find the number of integer solutions of
\[
    x_1+x_2+x_3+x_4=15
\]
satisfying
\[
    2\leq x_1\leq6,
    \qquad
    -2\leq x_2\leq1,
    \qquad
    0\leq x_3\leq6,
    \qquad
    3\leq x_4\leq8.
\]
\textit{Hint.} Use the inclusion-exclusion principle.
\end{problem}

\begin{solution}
Set
\[
    y_1=x_1-2,
    \qquad
    y_2=x_2+2,
    \qquad
    y_3=x_3,
    \qquad
    y_4=x_4-3.
\]
Then
\[
    y_1+y_2+y_3+y_4=12,
\]
with
\[
    0\leq y_1\leq4,
    \qquad
    0\leq y_2\leq3,
    \qquad
    0\leq y_3\leq6,
    \qquad
    0\leq y_4\leq5.
\]
Without the upper bounds, stars and bars gives
\[
    \binom{12+4-1}{4-1}=\binom{15}{3}
\]
nonnegative solutions.

Let $A_i$ be the set of solutions violating the upper bound on $y_i$. The single intersections have sizes
\[
    |A_1|=\binom{10}{3},
    \quad
    |A_2|=\binom{11}{3},
    \quad
    |A_3|=\binom{8}{3},
    \quad
    |A_4|=\binom{9}{3}.
\]
The nonzero pairwise intersections have sizes
\[
\begin{aligned}
    |A_1\cap A_2|&=\binom{6}{3},&
    |A_1\cap A_3|&=\binom{3}{3},&
    |A_1\cap A_4|&=\binom{4}{3},\\
    |A_2\cap A_3|&=\binom{4}{3},&
    |A_2\cap A_4|&=\binom{5}{3}.&&
\end{aligned}
\]
All other pairwise intersections and all intersections of three or more of the $A_i$ are empty. Therefore, by inclusion-exclusion, the required number is
\[
\begin{aligned}
    \binom{15}{3}
    &-\left(
        \binom{10}{3}+\binom{11}{3}+\binom{8}{3}+\binom{9}{3}
      \right)\\
    &+\left(
        \binom{6}{3}+\binom{3}{3}+\binom{4}{3}
        +\binom{4}{3}+\binom{5}{3}
      \right)\\
    &=69.
\end{aligned}
\]
\end{solution}

\begin{problem}{4}
Let $n\in\N$ and $k\in\Z$ with $0\leq k\leq n$. Prove that the number of $k$-element subsets of $[n]$ containing no two consecutive integers is
\[
    \binom{n-k+1}{k}.
\]
\end{problem}

\begin{solution}
The case $k=0$ is immediate. Assume $k\geq1$, and write such a subset as
\[
    \{a_1<a_2<\cdots<a_k\}.
\]
Since no two selected integers are consecutive,
\[
    a_{i+1}\geq a_i+2
\]
for every $1\leq i<k$. Define
\[
    b_i=a_i-(i-1),
    \qquad 1\leq i\leq k.
\]
Then
\[
    1\leq b_1<b_2<\cdots<b_k\leq n-k+1,
\]
so $\{b_1,\ldots,b_k\}$ is a $k$-element subset of $[n-k+1]$.

Conversely, from any
\[
    1\leq b_1<b_2<\cdots<b_k\leq n-k+1,
\]
defining $a_i=b_i+(i-1)$ gives a $k$-element subset of $[n]$ with no two consecutive elements. These constructions are inverse to each other. Hence the desired subsets are in bijection with the $k$-element subsets of $[n-k+1]$, and their number is
\[
    \binom{n-k+1}{k}.
\]
\end{solution}

\begin{remark}
The cyclic analogue is commonly known as \emph{Cayley's problem}: one selects $k$ vertices of a regular $n$-gon with no two selected vertices consecutive.
\end{remark}

\begin{problem}{5}
Let $n\in\N$. An $n$-derangement is a permutation $\pi\in S_n$ with no fixed point, i.e.
\[
    \pi(i)\neq i
    \qquad\text{for every }i\in[n].
\]
Let $D_n$ denote the number of $n$-derangements.
\begin{problemsubparts}
    \item Prove that
    \[
        D_n=n!\sum_{k=0}^{n}\frac{(-1)^k}{k!}.
    \]

    \item Prove that, for every $n\geq2$,
    \[
        D_n=nD_{n-1}+(-1)^n.
    \]
\end{problemsubparts}
\end{problem}

\begin{solution}
\begin{problemsubparts}
    \item For each $i\in[n]$, let
    \[
        A_i=\{\pi\in S_n\mid \pi(i)=i\}.
    \]
    The permutations that are not derangements are precisely the elements of $A_1\cup\cdots\cup A_n$. For any distinct indices $i_1,\ldots,i_k$,
    \[
        |A_{i_1}\cap\cdots\cap A_{i_k}|=(n-k)!,
    \]
    since those $k$ points are fixed and the remaining $n-k$ elements may be permuted arbitrarily. Hence inclusion-exclusion gives
    \[
    \begin{aligned}
        |A_1\cup\cdots\cup A_n|
        &=\sum_{k=1}^{n}(-1)^{k+1}\binom{n}{k}(n-k)!\\
        &=n!\sum_{k=1}^{n}\frac{(-1)^{k+1}}{k!}.
    \end{aligned}
    \]
    Therefore
    \[
    \begin{aligned}
        D_n
        &=n!-|A_1\cup\cdots\cup A_n|\\
        &=n!\sum_{k=0}^{n}\frac{(-1)^k}{k!}.
    \end{aligned}
    \]

    \item Using part (a),
    \[
    \begin{aligned}
        nD_{n-1}
        &=n!\sum_{k=0}^{n-1}\frac{(-1)^k}{k!}\\
        &=n!\left(
            \sum_{k=0}^{n}\frac{(-1)^k}{k!}
            -\frac{(-1)^n}{n!}
          \right)\\
        &=D_n-(-1)^n.
    \end{aligned}
    \]
    Rearranging yields
    \[
        D_n=nD_{n-1}+(-1)^n.
    \]
\end{problemsubparts}
\end{solution}


\end{document}
